\chapter*{Abstract}
\label{sec:Abstract}

\noindent The current dynamics of the labor market, driven by automation and artificial intelligence, demand precise tools for the standardization of competencies and occupational classifications. In this context, the automatic matching of job posting titles presents itself as a fundamental challenge in natural language processing.

This paper describes the system developed for participation in the task of \glslink{JTM}{Job Title Matching} of the \glslink{Tarea}{TalentCLEF} competition, which aims to address the case for the Spanish domain. A cascaded modular architecture is proposed, designed with the intention of maximizing accuracy in candidate retrieval. The system integrates three sequential stages: a dense retrieval phase based on three bi-encoder models optimized through supervised learning and contrastive loss functions; a rank fusion stage that combines the predictions to improve robustness; and a final re-ranking phase using a cross-encoder model to refine candidate selection.

The experimental development included an exhaustive analysis of models in a \textit{zero-shot} scenario, revealing the superiority of multilingual architectures over monolingual ones, even for the Spanish language context. The obtained results validate the effectiveness of the proposed approach: the system reached third place among participants in the overall ranking, with a mean \gls{MAP} score of 0.521 across the three required submission languages. However, of the three languages (English, German, and Spanish), performance was lowest in the target language. The subsequent analysis discusses the discrepancy between validation and test performance, as well as the paradox of achieving better results in other languages despite specific optimization for Spanish, underscoring the importance of data diversity and the intrinsic biases that characterize massively pre-trained models.
